\section{Background and Problem}
\label{sec:background}
\paragraph{Policy semantics.}
NetworkPolicy isolates ingress and egress independently. Once any policy
selects a Pod in a direction, only the union of the applicable policies'
rules is permitted in that direction; rules add permissions and never
remove them~\cite{k8snetworkpolicy}. A connection between two isolated
Pods must pass the source's egress rules and the destination's ingress
rules, and reply traffic of an allowed connection is implicitly allowed.
The semantics are connection-oriented, which makes reusing an admission
decision for the rest of a connection natural; the specification
prescribes no mechanism.

\paragraph{Lifecycle versus synchronization.}
Policies and Pods are created asynchronously and the API exposes no
barrier that says a policy is enforced everywhere. The expectation is
narrower: once an agent has processed a policy, a new Pod it selects must
be isolated before its containers start; allow rules may converge later.
A system can therefore be safe while delaying legitimate traffic, and a
fast transition to \texttt{Running} proves neither safety nor
availability. Here a connection gate ties the two together
(Section~\ref{sec:evaluation}).

\paragraph{Four quantities.}
``Identity scale'' conflates four things: live endpoints; distinct
policy-relevant label sets; allocated numeric identities, which can
outlive the Pods that created them; and materialized datapath entries,
which depend on the representation. Several Pods share one identity, one
identity contributes an entry to every endpoint that selects it, and
permitted pairs are unrelated to active connections.
Section~\ref{sec:evaluation} says which quantity each experiment varies.

\paragraph{Where enforcement work is placed.}
An ordered chain costs work per rule visited, but ipset and nftables sets
are indexed and hold throughput across the set sizes Sutter
tested~\cite{sutter2017nftables}, and kube-proxy's nftables backend makes
update cost proportional to what changed~\cite{winship2025nftables}. NVP
and Open vSwitch translate logical networks into per-host actions and
cache flow decisions in the kernel with misses handled in
userspace~\cite{koponen2014nvp,pfaff2015ovs}; OVN's stateful ACLs use
\texttt{allow-related} backed by conntrack~\cite{ovnacl}. These are the
closest precedents for our design, and none is stateless. eBPF maps let
userspace mutate keyed state without replacing the
program~\cite{linuxbpfmaps}; Cilium keys policy on label-derived numeric
identities~\cite{ciliumidentities}, and the size of that identity space
and the capacity of a per-endpoint policy map are different limits. The
per-peer cross-product is not the only kernel-resident representation
either: OVS conjunctive match installs an $n\times m$ rule as $n+m$
flows~\cite{ovsfields}, Antrea compiles NetworkPolicy onto it and marks
established connections in conntrack~\cite{antreapipeline}, and Calico
factors selectors into ipsets referenced by a bounded number of
rules~\cite{calicoarch}.

\paragraph{Problem.}
Can an engine keep the kernel forwarding path for the bulk of traffic
while avoiding eager expansion of every permitted peer at every endpoint?
If so, where does the displaced work go, how does the admission point
obtain metadata fresh enough to be useful, and what happens to cached
decisions when policy changes? Packet throughput answers none of these.

\section{Design}
\label{sec:design}
\subsection{Connection Admission and Cached Decisions}
\sys{} splits enforcement into steering, done by nftables, and
evaluation, done by a userspace agent. The agent watches NetworkPolicies,
Namespaces, and Pods through the API and learns local Pod addresses from
the runtime. From the Pods it derives the \emph{managed} addresses, local
non-hostNetwork Pods selected by some policy, and keeps them in nftables
sets. A filter chain on the \texttt{postrouting} hook at priority
$\mathrm{NATSource}-5$, so that it runs before SNAT and sees IPVS traffic
that skips \texttt{forward}, accepts packets from local sockets owned by
UID~0, accepts packets whose connection carries the accept label
\emph{and} is in established or related state, and queues to NFQUEUE any
remaining packet whose source or destination is in a managed set. An
\texttt{input} chain does the same for Pod-to-host traffic, queueing only
\texttt{ct state new} (Figure~\ref{fig:architecture}).

\begin{figure}[t]
\centering
\begin{tikzpicture}[
	box/.style={draw,align=center,text width=2.65cm,minimum height=0.85cm,font=\small},
	flow/.style={-{Latex},thick},font=\scriptsize]
\node[box] (runtime) at (0,0) {Runtime / NRI\\local Pod addresses};
\node[box] (metadata) at (3.85,0) {Agent view: Pods,\\Namespaces, policies};
\node[box] (kernel) at (0,-1.65) {nftables steering,\\managed sets, conntrack};
\node[box] (evaluate) at (3.85,-1.65) {Userspace policy\\evaluation};
\node[box] (forward) at (0,-3.3) {Kernel forwarding\\for labeled connections};
\draw[flow] (runtime) -- (metadata);
\draw[flow] (metadata) -- (evaluate);
\draw[flow] (kernel.east) -- node[above]{NFQUEUE, uncached} (evaluate.west);
\draw[flow] (evaluate.south) -- ++(0,-0.45) -| node[pos=0.25,below]{verdict + ct label} (kernel.south east);
\draw[flow] (kernel.south) -- node[left,align=right]{accept label and\\established/related} (forward.north);
\end{tikzpicture}
\caption{Enforcement decomposition. Protocol and host-traffic exceptions
are omitted. API watches feed the agent's view; NRI supplies local
observations only, not remote policy synchronization.}
\label{fig:architecture}
\end{figure}

The callback parses the headers, resolves source and destination to Pod
metadata, evaluates egress on the source and ingress on the destination,
and returns a verdict. An accept verdict carries a conntrack label; from
then on the bypass rule matches every packet of that connection in the
kernel. Denied packets are dropped and receive no label, so the next
attempt on that flow is queued again.

Two properties of this cache matter. It is connection state, not a table
of permitted peers: cached state is bounded by connections actually made,
$F_{\ell}$ in Section~\ref{sec:model}, not by the peers a policy could
admit. And ``first-packet evaluation'' is an approximation: conntrack
assigns established state only after seeing the reply direction, so a
handshake's initial packets, retransmissions before the reply, and
protocol-specific transitions can each reach userspace. The queued rate
exceeds the connection rate and includes every denied attempt.

nftables keeps a role. It holds the steering rules and the managed sets,
and a change in which local Pods are selected is a set update; it does
not execute label selectors. Avoiding the peer cross-product does not
remove local rule maintenance, and the evaluator's indexes over Pods,
Namespaces, and policies cost CPU and memory. Once metadata is present
the decision is local; nothing on the packet path allocates a distributed
identity. Remote Pod metadata must still reach the node.

\subsection{Runtime Metadata and Startup Isolation}
If the agent learns its own node's Pod addresses only from the API, the
chain is: the runtime creates the sandbox and the CNI assigns an address;
the kubelet writes Pod status; the API server fans it out; the agent's
watch delivers it. Containers may start, and lifecycle hooks may begin
waiting on the network, before the last step; issue 85966 records the
resulting failures~\cite{k8s85966}. \sys{}'s NRI resolver populates the
local IP-to-Pod cache in its \texttt{RunPodSandbox} callback, from the
addresses the runtime reports or by inspecting the sandbox network
namespace, so local address discovery no longer waits for the
kubelet--API-server--watch round trip.

Two limits apply. Knowing an address is not enforcing on it: the steering
path must cover the Pod before its containers run, which depends on set
updates and callback ordering, not on NRI alone. And the hook is local:
policies, namespace labels, and remote peers still arrive asynchronously.
We claim removal of one asynchronous dependency, not zero startup
latency, and no NRI on/off comparison exists to quantify even that.

\subsection{Policy Changes and Failure Semantics}
Caching makes policy updates part of the correctness argument. The
specification leaves the effect of an update on established connections
implementation-defined~\cite{k8snetworkpolicy}. \sys{}'s strict mode
scans conntrack, reevaluates supported established flows, and clears the
accept label where current policy no longer permits the flow, so the next
packet is queued and re-decided. Revocation latency is watch delivery
plus scan scheduling, scan duration, and packet arrival, and scan cost
grows with the flows inspected.

\paragraph{Queue overflow.}
Any design that decides in userspace behind a bounded kernel queue must
choose what the kernel does when the queue is full. NFQUEUE offers two
answers, drop or accept unevaluated~\cite{nfqueueapi}; \sys{} exposes
the choice as a flag and its queue holds 1,024 packets. This is the
counterpart of the eager design's overflow choice, where Cilium can lock
down an endpoint whose policy map is full~\cite{ciliumagentref}: in both
places the operator trades availability against isolation at a capacity
boundary. Denied packets are never labeled, so denied traffic is part of
the queue's load rather than of the cached path. The archived runs did
not measure queue occupancy or drops and did not approach overload, so
they characterize neither mode; the single-node experiment in
Section~\ref{sec:evaluation} is where that boundary would be measured.

\paragraph{Unknown peers fail closed.}
When the evaluator cannot resolve a remote address to a Pod, the peer is
\texttt{nil}, and a \texttt{nil} peer matches only \texttt{ipBlock}
rules. Under default deny the packet is dropped. That is the correct
isolation behavior and an availability cost: a permitted connection from
a freshly created remote Pod is denied until the destination node has
that Pod's metadata. NRI shortens this only for local Pods; for remote
peers the node depends on watch delivery or on the distribution tier of
Section~\ref{sec:iptracker}. The benchmark exercises exactly this path: a
sandbox's first connection is admitted at the gateway node only after
that node has learned the sandbox's address, so the post-scheduling
latencies in Section~\ref{sec:evaluation} contain this convergence.

\paragraph{Safety property and trust boundary.}
The property we want is conditional: for a policy an agent has already
processed, a newly selected Pod must not send or receive forbidden
traffic before enforcement applies; a permitted connection may be delayed
by missing metadata, and that delay is not a violation. The archived runs
do not establish it. They exercise the permitted sandbox-to-gateway path
on every Pod and the forbidden path once per cluster, by hand,
unarchived; missing metadata, label changes, and address reuse are
untested. The
trust boundary includes the node kernel, the runtime, the agent, and the
CNI's binding of addresses to Pods. A workload that can spoof a peer's
address defeats the evaluator; hostNetwork, privileged Pods, and address
translation need explicit scope that NFQUEUE and NRI do not provide.

\section{Where the Scaling Cost Moves}
\label{sec:model}
Let $I$ be the distinct peer identities a policy selects, $P_{\ell}$ the
local endpoints it applies to, $d$ the directions materialized, and
$F_{\ell}$ the active connections tracked on the node. An engine that
stores one entry per peer identity per direction at each endpoint carries
\begin{equation}
E_{\mathrm{endpoint}} \approx dI,\qquad
E_{\mathrm{node}} \approx dIP_{\ell}.
\label{eq:materialization}
\end{equation}
This describes one representation, not a lower bound: sharing,
wildcards, aggregation, or a factored match (Section~\ref{sec:related})
change it, other rules add to it, and entries are not bytes. Under the
same assumptions a new identity selected by every local endpoint costs
about $dP_{\ell}$ writes, and a rate $\lambda_I$ of new identities costs
$dP_{\ell}\lambda_I$ writes per node before dissemination and selector
recomputation. Replacing Pods that reuse a label set leaves this nearly
unchanged; fresh-identity churn is a different workload from Pod churn.
The benchmark tiers vary $I$; whether their churn rounds create fresh
identities is not recorded.

Deferred evaluation removes the $dIP_{\ell}$ term and keeps everything
else:
\begin{equation}
M_{\mathrm{deferred}} = M_{\mathrm{metadata}} +
M_{\mathrm{steering}} + M_{\mathrm{conntrack}}(F_{\ell}).
\end{equation}
This is an accounting, not a proof of lower total memory: both designs
carry conntrack, and the userspace Pod and policy indexes can dominate.
What is removed is materialized permission state; what replaces it is
work at admission time.

That work has a capacity. Let $\lambda_q$ be the rate of packets reaching
the queue, including denied attempts and pre-establishment
retransmissions, and $\mu$ the rate at which the agent returns verdicts.
The first-order condition for the queue not to grow is
\begin{equation}
\rho = \lambda_q/\mu < 1,
\end{equation}
and because the queue is 1,024 packets deep, exceeding it is the
overflow boundary of Section~\ref{sec:design}, not only latency;
bursts and CPU contention can reach it below $\rho=1$. Our only
measurement of $\mu$ is the repository's three-node ApacheBench
run~\cite{knptesting}: 10,000 non-keepalive requests at concurrency 1,000
completed at 2,317 requests/s, P50 5~ms, P99 3,080~ms dominated by
connect time, with the kernel logging \texttt{nf\_conntrack: table full}
at the default 262,144 entries; its charts show queued-packet peaks near
8,000/s per node and callback durations of 80--120~$\mu$s median and
170--630~$\mu$s P99. It is a 2024 measurement with no hardware, kernel,
or agent version recorded, and conntrack exhaustion confounds the tail;
we treat it as an order of magnitude for one node, not as $\mu$. The
720-node runs measured neither $\lambda_q$ nor $\mu$.

The favorable regime is therefore precise: a permission graph that is
large or changing fast relative to the connections actually attempted,
with $\rho$ well below one on every node. A dense permission graph with
sparse traffic fits. A dense connection graph, frequent revocation,
sustained denied traffic, or many short connections does not; there the
eager engine pays once per identity and the deferred engine pays per
connection.

\section{Implementation}
\label{sec:implementation}
The inspected code is the Go kube-network-policies repository at revision
\texttt{9984dcd}. The dataplane controller programs nftables over netlink
and opens NFQUEUE with a maximum queue length of 1,024 packets, a copy
range of 128 bytes, GSO enabled, and a 100~ms write timeout. These are
configuration constants, not measured service rates.

The callback times each packet, records the duration in a histogram
labeled by protocol and family, and counts verdicts by protocol, family,
and outcome; queue occupancy and drops are exported as well. Callback
time starts at delivery to userspace, so it excludes queueing before
delivery and the verdict's return path; it is not connection-setup
latency as an application sees it. None of these series exist for the
archived runs: the deployed Kindnet build exposes no metrics port, and
every agent query in the harness returned no samples. DNS, ICMPv6
neighbor discovery, and loopback exceptions are additional rules omitted
from Figure~\ref{fig:architecture}.

The repository builds several agent flavors. The standard flavor queues
only traffic to or from managed addresses. The IPTracker flavor, at the
same revision, returns \emph{divert-all} from \texttt{ManagedIPs}: it
watches NetworkPolicy objects but has no Pod informer from which to
enumerate selected local Pods, so it queues every forwarded packet not
already accepted by conntrack, and the queue-load model applies to a
larger traffic fraction. The NRI resolver is an alternate source of local
Pod records, added and removed by runtime callbacks; namespace metadata
still comes from the API view.

One naming distinction is essential for reading the evaluation. \sys{}
names the engine inspected here; the benchmark labels complete network
stacks \emph{Kindnet} and \emph{Cilium}. The \sys{} cluster was created
with kops's \texttt{kindnet: \{\}} networking, patched to
\texttt{kindnet:v1.0.1} with an NRI socket mount by
\path{scripts/create-netpol.sh}; the repository's standalone install
manifest pins \texttt{kube-network-policies:v1.1.0}. Which policy
binary and flags were effective in each Kindnet run is not archived, and
neither image is matched to revision \texttt{9984dcd}. The latency
differences in Section~\ref{sec:evaluation} are differences between
stacks, not attributable solely to the code described here.
